\documentclass[]{article}

\usepackage[T1]{fontenc}
\usepackage[utf8]{inputenc}
\usepackage{amsthm, amsmath, amsfonts}
\usepackage{blkarray}
\usepackage{bigstrut}
\usepackage{graphicx}

\newcommand{\avaliacao}{3}
\newcommand{\disciplina}{Probabilidade e Estatística}
\newcommand{\semestre}{2024/1}


\theoremstyle{definition}
\newtheorem{ex}{Exercício}

\renewcommand{\SS}{\mathcal{S}}

\begin{document}
\thispagestyle{empty}
\begin{center}
\Large{\avaliacao ª Avaliação de \disciplina\ - \semestre}

\large{Prof. Dr. Alessandro José Queiroz Sarnaglia}
\end{center}

\begin{ex}
A resistência à ruptura de rebites de um fabricante possui valor médio de 7000 psi e variância de 160000 psi$^2$.
\begin{enumerate}
	\item Qual é a probabilidade da resistência à ruptura média de uma amostra aleatória de 49 rebites estar acima de 7050?% \textbf{Resp.:} 0,1908.
	\item Se o tamanho da amostra fosse 121, e não 49, qual seria a probabilidade do item anterior?% \textbf{Resp.:} 0,0846.
\end{enumerate}
\end{ex}


\begin{ex}
Seja $X_1, X_2, \ldots, X_n$ uma amostra aleatória de uma distribuição de Rayleigh com fdp
$$
f(x;\theta) = \frac{x}{\theta}\exp\left\{-\frac{x^2}{2\theta}\right\}, \ \ x > 0, \ \ \theta>0.
$$
\begin{enumerate}
	\item Encontre o $\hat\theta$, Estimador de Máxima Verossimilhança (EMV) de $\theta$;% \textbf{Resp.:} $\hat{\theta} = \frac{1}{2n}\sum_{i}X_i^2$;
	\item Encontre o vício $B(\hat{\theta})$ do EMV $\hat{\theta}$. \textit{Dica: Encontre $E(X_i^2)$, lembrando que cada $X_i$ também têm distribuição de Rayleigh e que a esperança da soma é a soma das esperanças.}% \textbf{Resp.:} Mostrando que $E(X_i^2) = 2\theta$, vem que $E(\hat{\theta}) = \frac{1}{2n}\sum_{i}E(X_i^2) = \frac{1}{2n}\sum_{i}2\theta = \frac{2n}{2n}\theta = \theta$. Logo $B(\hat\theta) = E(\hat{\theta}) - \theta = 0$.
	\item Encontre a estimativa de máxima verossimilhança de $\theta$ a partir das $n = 10$ observações a seguir:
	\begin{center}
		\begin{tabular}{ccccc}
			4,8 & 9,0 & 7,8 &14,7 &15,6\\
			27,3&21,2 &11,3 &11,3 &15,1\\
		\end{tabular}
	\end{center}
%	\textbf{Resp.:} $\hat\theta = $ 115,1225.
\end{enumerate}
\end{ex}


\begin{ex}
Um fabricante de aspirina afirma que os seus frascos de 25 unidades têm peso médio unitário (por comprimido) de 5 mg ou mais. Sendo possível assumir que o peso por comprimido segue distribuição normal, para testar a afirmação do fabricante, 25 comprimidos tirados de um lote bastante grande são pesados, resultando em um peso médio amostral unitário de 4,76 mg e um variância amostral unitária de 0,16 mg$^2$.
\begin{enumerate}
	\item Calcule o $p$-valor para testar as hipóteses apropriadas. %\textbf{Resp.:} 0,003.
	\item Utilizando o $p$-valor obtido acima e um nível de significância de 0,01, você rejeitaria $H_0$ ou não? Justifique. %\textbf{Resp.:} Rejeitaria, pois $p$-valor $= 0\text{,}003 < 0\text{,}01 = \alpha$.
\end{enumerate}
\end{ex}


\end{document}
